\section{Related Work}
\label{sec:related-work}

Research in IaC predominantly explores its adoption patterns, inherent challenges, quality assessment frameworks, and testing met-hodologies. Studies have proposed metrics and frameworks for evaluating IaC quality, often exemplified with tools like Ansible or through specific metrics tailored for Terraform artifacts \cite{konala2025frameworkmeasuringqualityinfrastructureascode, DALLAPALMA2020110726}. Furthermore, static analysis techniques for IaC are a critical research area, focusing on identifying security vulnerabilities, compliance deviations, and best practice violations pre-deployment \cite{9779848}. The evaluation of IaC tools typically encompasses assessments of their performance, availability, and the efficancy of automated validation processes in ensuring robust, error-free cloud deployments \cite{article, 10516612}. The convergence of AI and IaC is an emerging frontier, with ongoing research investigating how generative AI can automate DevOps configurations, scripting, and workflow management, particularly in generating Terraform and Unusable scripts directly from design specifications \cite{article_kate, kumar_ai-powered_2025}.


Existing literature extensively covers the practical implications of AI in software development, with GitHub Copilot, an AI-powered coding assistant developed by GitHub and OpenAI, being a primary subject of empirical investigation. Numerous studies have assessed its efficiency, impact on developer productivity, and the quality of its code suggestions \cite{pandey2024transformingsoftwaredevelopmentevaluating, 9796235}. Findings consistently indicate that tools like GitHub Copilot can significantly expedite coding tasks, reduce cognitive load, and enable developers to concentrate on higher-level problem-solving \cite{kalliamvakou_research_2022}. Nevertheless, evaluations also underscore persistent challenges concerning the correctness, security, and overall understandability of AI-generated code, emphasizing the indispensable role of rigorous testing and human oversight \cite{10.1145/3715108, oertel_dont_2025}. Comparative analyses frequently pit GitHub Copilot against other AI code generation tools, such as Amazon CodeWhisperer and ChatGPT , using various metrics including code quality, functional accuracy, and adherence to established coding standards \cite{yetistiren2023evaluating} . The performance of these tools is typically gauged by their ability to produce syntactically correct and functionally accurate code, alongside their influence on developer workflows and satisfaction \cite{oertel_dont_2025}. While their potential is undeniable, researchers continue to explore their limitations, particularly in addressing complex, context-dependent coding problems and ensuring the security of the generated output \cite{sun_terraform_2025}.

Despite the growing body of research on AI-assisted code generation and IaC, a notable gap exists concerning the comprehensive evaluation of newer, specialized tools within specific domains. While Windsurf is recognized as an AI-powered code editor, and its AI component Cascade is highlighted as an LLM-powered platform for real-time machine intelligence \cite{sewell_windsurf_2024, hofferber_my_2024, windsurf_editor, manuskhani_cascade_2024}, scientific literature specifically evaluating Windsurf Cascade in the context of Infrastructure as Code with Terraform remains limited. Discussions and reviews of Windsurf are primarily found in developer communities and blogs, lacking the rigorous academic scrutiny applied to more established tools like GitHub Copilot. This absence of peer-reviewed studies on Windsurf Cascade's performance, code quality, or specific applications in IaC with Terraform underscores a critical research void. Our study directly addresses this gap by providing a systematic evaluation of VSCode GitHub Copilot and Windsurf Cascade in the IaC domain with Terraform, assessing the quality, functionality, and completeness of their generated code across a predefined set of tasks. By doing so, we aim to contribute novel empirical data and insights into the comparative efficiency of these tools, thereby enriching the academic discourse on AI-assisted IaC development and offering practical guidance for practitioners.

In literature, the main focus is on methods combining quantitative and qualitative metrics to assess code quality, correctness, and efficiency. Key metrics include correctness and functional accuracy, verified by benchmarks; code quality, which considers readability, maintainability, and efficiency through measures such as cyclomatic complexity and static analysis warnings and security, evaluated using static analysis tools and manual reviews to detect vulnerabilities \cite{sabra2025assessingqualitysecurityaigenerated}. Productivity gains are assessed by comparing task completion times and cognitive load with and without AI assistance \cite{hendrich_research_2024}, while user perception and trust are studied through surveys and interviews \cite{10685197}. Empirical studies often involve experiments where AI-generated solutions to predefined tasks are compared against human or alternative AI code \cite{10825958}. Due to the non-deterministic nature of large language models, robust evaluation designs are necessary \cite{10.1145/3697010}. For domain-specific contexts like Infrastructure as Code with Terraform, evaluations must also address resource optimization, and cloud best practices. Our methodology aligns with these principles and adapts them to the specific requirements of Terraform-based IaC.